\section{Base algorithm}
\label{appendix:analysis of the main algorithm}

In this section we focus on our base algorithm and analysis: we design an algorithm that uses a distributional prediction in an essentially optimal way (as the lower bounds discussed earlier show).  We first give our base algorithm formally in Algorithm~\ref{alg:distributional truth algo}. This algorithm guarantees an upper bound (Theorem~\ref{thm:main}). We present the intuition of our proof (Appendix~\ref{sec:intuition of base proof}), followed by a full proof of it (Appendix~\ref{sec:full proof of base algorithm}). Without loss of generality, we assume that $\sqrt{b}$ is an integer; otherwise, we can use $\lfloor \sqrt{b} \rfloor$ instead. 


\begin{algorithm}
\caption{Base algorithm, with expected additive loss bounded by $O(\sqrt{b}\max(\mathrm{EMD},1))$}
\label{alg:distributional truth algo}
\begin{algorithmic}[1]
\IF{$A_{\hat{K}}$ is the optimal policy for $\hat{p}$ with some finite $\hat{K}$}
    \STATE Buy at the beginning of day $\hat{K}+\sqrt{b}+1$, i.e., $ALG = A_{\hat{K}+\sqrt{b}}$
\ELSE
    \STATE Keep renting, i.e., $ALG = A_{\infty}$
\ENDIF
\end{algorithmic}
\end{algorithm}

In other words, Algorithm~\ref{alg:distributional truth algo} distinguishes two cases depending on the optimal policy under the predicted distribution $\hat p$.  If the optimal policy for $\hat p$ is of the form $A_{\hat K}$ for some finite $\hat K$, the algorithm delays the purchase time by an additive $\sqrt b$ and buys at time $\hat K + \sqrt b + 1$. Otherwise, when the optimal policy for $\hat p$ is $A_{\infty}$, the algorithm simply keeps renting and never buys.


Recall the statement of Theorem~\ref{thm:main}:

\maintheorem*

To prove Theorem~\ref{thm:main},
we break into cases depending on the structure of the optimal policies for $\hat p$ and $p$. We first prove Theorem~\ref{thm:main} when the optimal policy for $\hat p$ is $A_{\infty}$ (Lemma~\ref{lm:phat Ainf and p AK with EMD O1}). When the optimal policy for $\hat{p}$ is $A_K$ with finite $K$ the analysis is more complicated, and we need to break into two further cases depending on the true optimal policy for $p$. In Lemma~\ref{lm:phat is AK and p is Ainf and root b loss} we prove the theorem for the optimal policy (for $p$) being $A_{\infty}$ and in Lemma~\ref{lm:phat is AKhat and p is AK root b loss} we consider the case when both the optimal policies for $p$ and for $\hat p$ involve buying at some finite time.  Together, these lemmas directly imply Theorem~\ref{thm:main}.  


\subsection{Intuition of proof}
\label{sec:intuition of base proof}

While the full proof is in Appendix~\ref{sec:full proof of base algorithm}, we now provide some intuition about it.
First, it is helpful for us to break into cases depending on the optimal policies for $p$ and $\hat p$ (as discussed above) because once we know the optimal policy we can explicitly write down the cost of the optimal policy, which allows us to write $\text{diff}_p$ explicitly.  

Now suppose that the true distribution were actually $\hat p$ rather than $p$.  Then clearly our algorithm is optimal if this policy is $A_{\infty}$, and otherwise is suboptimal.  But we can also consider what happens if we run $OPT_p$ (the optimal policy for $p$) when the truth is $\hat p$, and we can define $\text{diff}_{\hat p}$ as the expected cost of our algorithm minus the expected cost of $OPT_p$ when the true distribution is $\hat p$.  So $\text{diff}_{\hat p}$ is nonpositive when our policy is $A_{\infty}$ (since in this case we are optimal), but otherwise can be quite large.

If it \emph{were} true that $\text{diff}_{\hat p}$ was always small, then a natural strategy would be to start with that as a baseline and analyze how the $\text{diff}_p$ increases when compared with $\text{diff}_{\hat p}$.  In other words, we could try to argue that $\text{diff}_p - \text{diff}_{\hat p} \leq O(\sqrt{b}\max(EMD(\hat{p},p),1))$, which would then imply Theorem~\ref{thm:main} if $\text{diff}_{\hat p}$ were small enough.  Unfortunately this runs into two related issues.  First, $\text{diff}_{\hat p}$ could be extremely negative, making our desired claim false.  Second, $\text{diff}_{\hat p}$ could be extremely large and positive, making the claim true but insufficient to prove Theorem~\ref{thm:main}.  Fortunately, we can get around both of these issues in the same way: we can simply normalize by $\text{diff}_{\hat p}$.  More formally, we can prove that $\text{diff}_p - \text{diff}_{\hat{p}} \leq O(\sqrt{b}\max(EMD(\hat{p},p),1)) - \text{diff}_{\hat{p}}$.  Simply adding $\text{diff}_{\hat p}$ to both sides then implies the theorem.  

So our strategy is try to bound $\text{diff}_p - \text{diff}_{\hat{p}}$.  This is about the change in the $\text{diff}$ function when the distribution changes from $\hat p$ to $p$, and so a natural approach is to utilize the optimal transport plan from $\hat p$ to $p$.  We call each value $f(x,y)$ in this transport plan a \emph{map}, since it tells us how much probability mass to move from $x$ to $y$.  Some of the maps either benefit us (i.e., result in $\text{diff}_p$ decreasing compared to $\text{diff}_{\hat p}$) or result in an increase that can still be bounded in simple ways (e.g., directly bounded by $EMD(\hat p, p)$). We call such maps \emph{good maps} and all others are \emph{bad maps}, and only need concern ourselves with bad maps.  Although which maps are bad depends on which precise case we are in, we generally split the bad maps into two types: maps that move mass from $t_1$ to $t_2$ with $|t_1-t_2| \geq \sqrt{b}$ and maps that move mass from $t_1$ to $t_2$ with $|t_1-t_2| < \sqrt{b}$. If a map is in the regime where the moving distance is large (larger than $\sqrt{b}$), the amount of mass moved by this map is then bounded by $\frac{EMD(\hat p, p)}{\sqrt{b}}$ by the definition of $EMD(\hat p, p)$. This is a small enough amount of mass that we can straightforwardly argue that the increase it causes lies within our desired bound.

If a map is in the regime where the moving distance is small (smaller than $\sqrt{b}$), then we instead directly use the fact that our policy is related to (though not identical to) the optimal policy for $\hat p$.  In particular, an obvious observation is that when moving mass from $\hat{p}$ to $p$ using those bad maps in the optimal transport plan, the amount of mass one can move from point $t$ is at most $\hat{p}_t$. So if we can argue that $\hat p_t$ is small, then we get that these types of bad maps move a small amount of mass a small distance, and so the increase they cause in the $\text{diff}$ function still lives within our desired bound.  And it turns out that we can indeed show this, although the precise guarantee and strategy is different for different cases. 
 But for example, when the optimal policy for $\hat{p}$ is $A_K$, it is intuitively true (and can be made formal) that $\sum_{t=K}^{K+\delta} \hat p_t$ has to be small when $\delta$ is small (otherwise renting $K$ times then buying would not be the optimal policy; it would be better for us to wait a little longer to buy). Therefore, we leverage the mass distribution in 
$\hat{p}$ to connect the change in additive loss $\text{diff}_p - \text{diff}_{\hat{p}}$ with the required bound $O(\sqrt{b}\max(EMD(\hat{p},p),1)) - \text{diff}_{\hat{p}}$. 


\subsection{Full proof}
\label{sec:full proof of base algorithm}

This section gives the technical proofs of Theorem~\ref{thm:main}.
we break into cases depending on the optimal policies for $\hat{p}$ and $p$ to get the claimed bound.


\begin{lemma}
    \label{lm:phat Ainf and p AK with EMD O1}
    Let $\hat{p}$ be the prediction and $p$ be the unknown truth. If the optimal policy for $\hat{p}$ is $A_\infty$, then taking $ALG=A_\infty$, we have $\text{diff}_p = c(ALG_p) - c(OPT_p) \leq O(\sqrt{b} \max(EMD(\hat{p},p),1))$.
\end{lemma}

\begin{proof}
    If the optimal policy for $p$ is also $A_\infty$, then $\text{diff}_p = 0$ and the required bound holds trivially. Thus, we only need to show that the bound holds when the optimal policy for $p$ is $A_K$ for some finite $K$.

    First, we have 
    $$\text{diff}_p = \sum_{t=1}^N tp_t - (\sum_{t \leq K}tp_t + \sum_{t>K}(K+b)p_t) = \sum_{t>K}(t-(K+b))p_t.$$ 
    Since $\text{diff}_p$ is a function of $p$, we similarly define $\text{diff}_{\hat{p}}$ as the quantity by using the same policies but the different distribution $\hat{p}$ rather than $p$. In other words, $\text{diff}_{\hat{p}} = c(ALG_{\hat{p}}(\hat{p})) - c((OPT_p)_{\hat{p}})$. Then, we consider the optimal transport plan from $\hat{p}$ to $p$ and examine the change of the additive loss $\text{diff}_p - \text{diff}_{\hat{p}}$. We want to show that $\text{diff}_p - \text{diff}_{\hat{p}} \leq O(\sqrt{b}\max(EMD(\hat{p},p),1)) - \text{diff}_{\hat{p}}$. That is, we want this change in additive loss to be bounded by $O(\sqrt{b}\max(EMD(\hat{p},p),1)) - \text{diff}_{\hat{p}}$.

    Note that $A_\infty$ is the optimal policy for $\hat{p}$, we have $\text{diff}_{\hat{p}} \leq 0$ and 
    \[\text{diff}_{\hat{p}} = \sum_{t>K}(t-(K+b))\hat{p}_t = \sum_{t>K+b}(t-(K+b))\hat{p}_t - \sum_{K<t<K+b}(K+b-t)\hat{p}_t \leq 0.\]
    We divide the domain into 3 regions: $[1,K], (K,K+b], (K+b,\infty)$. There are two kinds of maps in the optimal transport plan: from some point in $(K,K+b]$ to some point in $[1,K]$, or not in such a form. If we consider how a map from one region to any region (including itself) affects the change in additive loss, notice that for any map that is not from $(K,K+b]$ to $[1,K]$, the change in additive loss by the map is at most its own contribution to the Earth Mover's Distance.
    To argue this, if a map is from~$[1,K]$ to~$[1,K]$, then it does not affect the change in additive loss. If a map is from~$[1,K]$ to~$(K,K+b]$, then it only decreases the value and hence benefits us. If a map is from~$t_1 \in [1,K]$ to~$t_2 \in (K+b,\infty)$ with mass~$\epsilon$, then it increases the value by~$(t_2 - (K+b))\epsilon < (t_2-t_1)\epsilon$ which equals its contribution to the~$EMD$. If a map is from~$t_1 \in (K,K+b]$ to~$t_2 \in (K,K+b]$ with mass~$\epsilon$, then it changes the value by~$(K+b-t_1)\epsilon - (K+b-t_2)\epsilon \leq |t_1-t_2|\epsilon$ which equals its contribution to the~$EMD$. If a map is from~$t_1 \in (K,K+b]$ to~$t_2 \in (K+b,\infty)$ with mass~$\epsilon$, then it increases the value by~$(K+b-t_1)\epsilon + (t_2 - (K+b))\epsilon = (t_2-t_1)\epsilon$ which equals its contribution to the~$EMD$. If a map is from~$(K+b,\infty)$ to~$[1,K]$ or from~$(K+b,\infty)$ to~$(K,K+b]$, then again it only decreases the value and benefits us. If a map is from~$t_1 \in (K+b,\infty)$ to~$t_2 \in (K+b,\infty)$ with mass~$\epsilon$, then it changes the value by~$(t_2-(K+b))\epsilon - (t_1-(K+b))\epsilon = (t_2-t_1)\epsilon$ which equals its contribution to the~$EMD$.
    Thus, if we consider all such maps (not from~$(K,K+b]$ to~$[1,K]$) as a group, it will increase the additive loss by some amount bounded by $EMD(\hat{p},p)$ which is at most a constant. We call such maps in this easy case the good maps and those we need to take further actions the bad maps.

    For maps (elements in the optimal transport plan) from $(K,K+b]$ to $[1,K]$, we have $\text{diff}_p - \text{diff}_{\hat{p}} = \sum_{K<t<K+b}(K+b-t)(\hat{p}_t - p_t)$. Let us first consider the set of bad maps $\mathcal{B}_1$ from $(K+\sqrt{b},K+b]$ to $[1,K]$. Since the moving distance for such map is at least $\sqrt{b}$, the total mass moving by such bad maps is at most $O(\frac{EMD}{\sqrt{b}})$. Thus, 
    $$\left. \text{diff}_p - \text{diff}_{\hat{p}} \right|_{\mathcal{B}_1} \leq (K+b-(K+\sqrt{b}))\sum_{K+\sqrt{b}<t\leq K+b}\left. \epsilon_t \right|_{\mathcal{B}_1} \leq (b-\sqrt{b})\cdot O(\frac{EMD}{\sqrt{b}})=O(\sqrt{b}\cdot EMD),$$
    where $\left. \epsilon_t \right|_{\mathcal{B}_1} = \left. \hat{p}_t - p_t \right|_{\mathcal{B}_1}$ is the mass change by all such bad maps in $\mathcal{B}_1$ at $t$ for any point $K+\sqrt{b}<t \leq K+b$.

    Then, let us consider the set of bad maps $\mathcal{B}_2$ from $(K,K+\sqrt{b}]$ to $[1,K]$. We will discuss the change in additive loss by such bad maps in $\mathcal{B}_2$. That is, we want to bound 
    $$\left. \text{diff}_p - \text{diff}_{\hat{p}} \right|_{\mathcal{B}_2} = \sum_{K<t \leq K+\sqrt{b}}(K+b-t) \left. (\hat{p}_t - p_t) \right|_{\mathcal{B}_2}= \sum_{K<t \leq K+\sqrt{b}}(K+b-t)\left. \epsilon_t \right|_{\mathcal{B}_2},$$
    where $\left. \epsilon_t \right|_{\mathcal{B}_2} = \left. \hat{p}_t - p_t \right|_{\mathcal{B}_2}$ is the mass change by all such bad maps in $\mathcal{B}_2$ at $t$ for any point $K<t \leq K+\sqrt{b}$. Note that it suffices to show that $\left. \text{diff}_p - \text{diff}_{\hat{p}} \right|_{\mathcal{B}_2} \leq O(\sqrt{b})- \text{diff}_{\hat{p}}$ to get our required result.

    Recall that in $\hat{p}$, $A_\infty$ is the optimal policy. Thus, we have
     \begin{align}
    & c((A_\infty)_{\hat{p}}) \leq c((A_{K+\sqrt{b}})_{\hat{p}}) \notag \\
    \iff & \sum_t t\hat{p}_t \leq \sum_{t \leq K+\sqrt{b}} t\hat{p}_t + \sum_{t>K+\sqrt{b}}(K+\sqrt{b}+b)\hat{p}_t \notag \\
    \iff & \sum_{t > K+\sqrt{b}} t\hat{p}_t \leq \sum_{t>K+\sqrt{b}}(K+\sqrt{b}+b)\hat{p}_t \notag \\
    \iff & \sum_{t > K+b} (t- (K+\sqrt{b}+b))\hat{p}_t \leq \sum_{K+\sqrt{b}<t \leq K+b}(K+\sqrt{b}+b-t)\hat{p}_t \notag \\
    \iff & \sum_{t > K+b} (t- (K+b))\hat{p}_t \leq \sum_{t > K+b}\sqrt{b}\hat{p}_t + \sum_{K+\sqrt{b}<t \leq K+b}(K+\sqrt{b}+b-t)\hat{p}_t. \quad \label{eq:dagger} \tag{\(\dagger\)}
    \end{align}

    By the definition of $\left. \epsilon_t \right|_{\mathcal{B}_2}$, we know that $\left. \epsilon_t \right|_{\mathcal{B}_2} \leq \hat{p}_t$ for $K<t \leq K+\sqrt{b}$. Therefore, to show that $\left. \text{diff}_p - \text{diff}_{\hat{p}} \right|_{\mathcal{B}_2} \leq O(\sqrt{b})- \text{diff}_{\hat{p}} = O(\sqrt{b})- (\sum_{t>K}(t-(K+b))\hat{p}_t)$, it suffices to show that
    \[\sum_{K<t \leq K+\sqrt{b}}(K+b-t)\hat{p}_t \leq O(\sqrt{b})- \sum_{t>K}(t-(K+b))\hat{p}_t\]
    which is equivalently
    \[\sum_{t>K+b} (t-(K+b))\hat{p}_t \leq O(\sqrt{b}) + \sum_{K+\sqrt{b}<t\leq K+b} (K+b-t)\hat{p}_t.\]
    The last line holds because of~\eqref{eq:dagger}.

    By combining results from good maps and bad maps, we finally get the required result.
\end{proof}

Next, we show that when the optimal policy for $\hat{p}$ is $A_K$ for some finite $K$, taking $ALG=A_{K+\sqrt{b}}$ implies the required bound. Before proving this, let us make some observations.

\begin{observation}
    \label{obs:AK plus root b is not too bad than AK}
    $c((A_{K+\sqrt{b}})_{\hat{p}}) \leq c((A_{K})_{\hat{p}}) + \sqrt{b}$.
\end{observation}

\begin{proof}
    We write down the definition of $c((A_{K+\sqrt{b}})_{\hat{p}})$ and $c((A_{K})_{\hat{p}})$. It is not hard to see that by taking the difference we have $c((A_{K+\sqrt{b}})_{\hat{p}}) - c((A_{K})_{\hat{p}}) \leq \sqrt{b} \sum_{t>K+\sqrt{b}}\hat{p}_t \leq \sqrt{b}$.
\end{proof}

\begin{claim}
    \label{claim:if AK is opt policy then K to K+T cannot have too much mass}
    Let $A_K$ be the optimal policy for $\hat{p}$. Then, for every $1 \leq T \leq b$, we have
    \[\sum_{K<t\leq K+T}\hat{p}_t \leq \frac{T}{b}\sum_{t>K}\hat{p}_t.\]
\end{claim}

\begin{proof}
    For $1 \leq T \leq b$, we will compare the cost of $A_K$ and $A_{K+T}$.
    \begin{align*}
    & c((A_K)_{\hat{p}}) \leq c((A_{K+T})_{\hat{p}}) \\
    \iff & \sum_{t \leq K} t\hat{p}_t + \sum_{t>K}(K+b)\hat{p}_t \leq \sum_{t \leq K+T} t\hat{p}_t + \sum_{t>K+T}(K+T+b)\hat{p}_t \\
    \iff & \sum_{K<t \leq K+T}(K+b)\hat{p}_t \leq \sum_{K<t \leq K+T} t\hat{p}_t + \sum_{t> K+T}T\hat{p}_t \\
    \iff & \sum_{K<t \leq K+T}(K+b-t)\hat{p}_t \leq  \sum_{t> K+T}T\hat{p}_t.
    \end{align*}
    
    Thus, we have 
    \begin{align*}
    b \sum_{K<t\leq K+T}\hat{p}_t \leq&  T \sum_{K<t \leq K+T} \hat{p}_t + \sum_{K<t \leq K+T}(K+b-t)\hat{p}_t\\
    \leq& T \sum_{K<t \leq K+T} \hat{p}_t + T \sum_{t> K+T} \hat{p}_t=T \sum_{t>K}\hat{p}_t.
    \end{align*}
    That is, $\sum_{K<t\leq K+T}\hat{p}_t \leq \frac{T}{b}\sum_{t>K}\hat{p}_t$.
\end{proof}

Now, we can start with the proof. In the following lemmas, we will analyze the cases separately based on the optimal policy for $p$ to be either $A_\infty$ or some $A_L$, given that the optimal policy for $\hat{p}$ is $A_K$ for some finite $K$ as well as $ALG=A_{K+\sqrt{b}}$ always.


\begin{lemma}
    \label{lm:phat is AK and p is Ainf and root b loss}
    Let $\hat{p}$ be the prediction and $p$ be the unknown truth. Assume that the optimal policy for~$\hat{p}$ is $A_K$ for some finite $K$. Let $ALG = A_{K+\sqrt{b}}$. If the optimal policy for $p$ is $A_\infty$, we have $\text{diff}_p = c(ALG_p) - c(OPT_p) \leq O(\sqrt{b} \max(EMD(\hat{p},p),1))$.
\end{lemma}

\begin{proof}
    We want to show that $\text{diff}_p = c((A_{K+\sqrt{b}})_p) - c((A_\infty)_p) \leq O(\sqrt{b} \max(EMD(\hat{p},p),1))$. Consider $\text{diff}_{\hat{p}} = c((A_{K+\sqrt{b}})_{\hat{p}}) - c((A_\infty)_{\hat{p}})$. Since $A_K$ is the optimal policy for $\hat{p}$, we have $c((A_{K})_{\hat{p}}) - c((A_\infty)_{\hat{p}}) \leq 0$. By Observation~\ref{obs:AK plus root b is not too bad than AK}, we have 
    \[\text{diff}_{\hat{p}} = c((A_{K+\sqrt{b}})_{\hat{p}}) - c((A_\infty)_{\hat{p}}) \leq c((A_{K})_{\hat{p}}) + \sqrt{b}- c((A_\infty)_{\hat{p}}) \leq \sqrt{b}.\]
    It turns out that if we can show $\text{diff}_p - \text{diff}_{\hat{p}} \leq O(\sqrt{b} \max(EMD(\hat{p},p),1))$, then $\text{diff}_p - \text{diff}_{\hat{p}} \leq O(\sqrt{b} \max(EMD(\hat{p},p),1)) - \text{diff}_{\hat{p}}$ and we are done.

    To show that $\text{diff}_p - \text{diff}_{\hat{p}} \leq O(\sqrt{b} \max(EMD(\hat{p},p),1))$, again we use the optimal transport plan and analyze the bad maps. After simple calculations, we have 
    $$\text{diff}_{\hat{p}} = \sum_{K+\sqrt{b}< t \leq K+\sqrt{b}+b}(K+\sqrt{b}+b-t)\hat{p}_t - \sum_{t > K+\sqrt{b}+b}(t-(K+\sqrt{b}+b))\hat{p}_t.$$ 
    Then, by considering maps between the regions $[1,K+\sqrt{b}]$, $(K+\sqrt{b},K+\sqrt{b}+b]$ and $(K+\sqrt{b}+b,\infty)$, the bad maps are those from $[1,K+\sqrt{b}]$ to $(K+\sqrt{b},K+\sqrt{b}+b]$ (other maps are again digested by their contributions to the $EMD$, as we did in Lemma~\ref{lm:phat Ainf and p AK with EMD O1}). 

    For maps $\mathcal{B}$ from $[1,K+\sqrt{b}]$ to $(K+\sqrt{b},K+\sqrt{b}+b]$ (elements in the optimal transport plan), since $\mathcal{B}$ has nothing to do with $t>K+\sqrt{b}+b$, we have $$\left. \text{diff}_p - \text{diff}_{\hat{p}}\right|_{\mathcal{B}} = \sum_{K+\sqrt{b}<t\leq K+\sqrt{b}+b}(K+\sqrt{b}+b-t)(p_t - \hat{p}_t).$$
    
    Let us first consider the set of bad maps $\mathcal{B}_1$ from $[1,K]$ to $(K+\sqrt{b},K+\sqrt{b}+b]$. Since the moving distance for such map is at least $\sqrt{b}$, the total mass moving by such bad maps is at most $O(\frac{EMD}{\sqrt{b}})$. Thus, 
    $$\left. \text{diff}_p - \text{diff}_{\hat{p}} \right|_{\mathcal{B}_1} = \sum_{ K+\sqrt{b}<t \leq K+\sqrt{b}+b}(K+\sqrt{b}+b-t)\left. \epsilon_t \right|_{\mathcal{B}_1} \leq b \cdot \left. \sum_{ K+\sqrt{b}<t \leq K+\sqrt{b}+b} \epsilon_t \right|_{\mathcal{B}_1} \leq O(\sqrt{b} \cdot EMD),$$
    where $\left. \epsilon_t \right|_{\mathcal{B}_1} = \left. p_t - \hat{p}_t \right|_{\mathcal{B}_1}$ is the mass change by all such bad maps in $\mathcal{B}_1$ at $t$ for any point $K+\sqrt{b}<t \leq K+\sqrt{b}+b$.

    Then, let us consider the set of bad maps $\mathcal{B}_2$ from $(K,K+\sqrt{b}]$ to $(K+\sqrt{b},K+\sqrt{b}+b]$. We similarly define $\left. \epsilon_t \right|_{\mathcal{B}_2} = \left. p_t - \hat{p}_t \right|_{\mathcal{B}_2}$ as the mass change by all such bad maps in $\mathcal{B}_2$ at~$t$ for any point $K+\sqrt{b}<t \leq K+\sqrt{b}+b$.
    Since we are considering bad maps of the form $\mathcal{B}_2$, we know that $\left. \sum_{ K+\sqrt{b}<t \leq K+\sqrt{b}+b} \epsilon_t \right|_{\mathcal{B}_2}$ is at most $\sum_{K<t \leq K+\sqrt{b}}\hat{p}_t$. Therefore, 
    $$\left. \text{diff}_p - \text{diff}_{\hat{p}} \right|_{\mathcal{B}_2} = \sum_{ K+\sqrt{b}<t \leq K+\sqrt{b}+b}(K+\sqrt{b}+b-t)\left. \epsilon_t \right|_{\mathcal{B}_2} \leq b \left. \sum_{ K+\sqrt{b}<t \leq K+\sqrt{b}+b} \epsilon_t \right|_{\mathcal{B}_2} \leq b \sum_{K<t \leq K+\sqrt{b}}\hat{p}_t.$$
    Again, this is bounded by $O(\sqrt{b})$ by taking $T=\sqrt{b}$ in Claim~\ref{claim:if AK is opt policy then K to K+T cannot have too much mass}.

    Therefore, we finish the proof as required.
\end{proof}

Now, the only left case is when $A_{\hat{K}}$ is the optimal policy for $\hat{p}$ and $A_K$ is the optimal policy for $p$ for some finite $\hat{K}$ and $K$. In this case, we refer to the optimal policy for $\hat{p}$ as $A_{\hat{K}}$ to avoid any ambiguity. We need to show that by taking $ALG = A_{\hat{K}+\sqrt{b}}$, we have $\text{diff}_p \leq O(\sqrt{b} \max(EMD(\hat{p},p),1))$. This needs a long proof, as one will see that we have to use a combination of techniques in the previous lemmas.

\begin{lemma}
    \label{lm:phat is AKhat and p is AK root b loss}
    Let $\hat{p}$ be the prediction and $p$ be the truth. Assume that the optimal policy for~$\hat{p}$ is $A_{\hat{K}}$ for some finite $\hat{K}$. Let $ALG = A_{\hat{K}+\sqrt{b}}$. If the optimal policy for $p$ is $A_K$ for some finite $K$, we have $\text{diff}_p = c(ALG_p) - c(OPT_p) \leq O(\sqrt{b} \max(EMD(\hat{p},p),1))$.
\end{lemma}

\begin{proof}
    We want to show that $\text{diff}_p = c((A_{\hat{K}+\sqrt{b}})_p) - c((A_K)_p) \leq O(\sqrt{b})$. Consider $\text{diff}_{\hat{p}} = c((A_{\hat{K}+\sqrt{b}})_{\hat{p}}) - c((A_K)_{\hat{p}})$. Since $A_{\hat{K}}$ is the optimal policy for $\hat{p}$, we have $c((A_{\hat{K}})_{\hat{p}}) - c((A_K)_{\hat{p}}) \leq 0$. By Observation~\ref{obs:AK plus root b is not too bad than AK}, we have 
    \[\text{diff}_{\hat{p}} = c((A_{\hat{K}+\sqrt{b}})_{\hat{p}}) - c((A_K)_{\hat{p}}) \leq c((A_{\hat{K}})_{\hat{p}}) + \sqrt{b}- c((A_K)_{\hat{p}}) \leq \sqrt{b}.\]
    Now, we want to find the bad maps in an optimal transport plan. We write down
    \[\text{diff}_{\hat{p}} = \left( \sum_{t \leq \hat{K}+\sqrt{b}} t \hat{p}_t + \sum_{t>\hat{K}+\sqrt{b}}(\hat{K}+\sqrt{b}+b)\hat{p}_t \right) - \left( \sum_{t \leq K} t \hat{p}_t + \sum_{t>K}(K+b)\hat{p}_t \right).\]
    We separate into two cases by comparing $\hat{K}+\sqrt{b}$ and $K$.

    \textbf{Case 1.} When $\hat{K}+\sqrt{b} < K$.

    In this case, by some simple calculations, we have 
    \begin{align*}
    \text{diff}_{\hat{p}} =& \sum_{\hat{K}+\sqrt{b}< t \leq K}(\hat{K}+\sqrt{b}+b-t)\hat{p}_t - \sum_{t > K}(K-(\hat{K}+\sqrt{b}))\hat{p}_t\\
    =& \sum_{t \leq \hat{K}+\sqrt{b}}(K-(\hat{K}+\sqrt{b}))\hat{p}_t + \sum_{\hat{K}+\sqrt{b}< t \leq K}(K+b-t)\hat{p}_t - (K-(\hat{K}+\sqrt{b})).
    \end{align*}
    By examining different kind of maps, one can check that there are two kinds of bad maps: either from $[1,\hat{K}+\sqrt{b}]$ to $(\hat{K}+\sqrt{b},K]$ (bad maps of the first kind $\mathcal{B}^1$), or from $(K,K+b]$ to $(\hat{K}+\sqrt{b},K]$ (bad maps of the second kind $\mathcal{B}^2$). We will analyze those bad maps of the first kind using the first expression of $\text{diff}_{\hat{p}}$ and analyze those bad maps of the second kind using the second expression in the above formula.

    For the set of bad maps of the first kind which are from $[1,\hat{K}]$ to $(\hat{K}+\sqrt{b},K]$, we denote it as $\mathcal{B}^1_1$. We know that maps in $\mathcal{B}^1_1$ has nothing to do with $t>K$. Moreover, since the moving distance for such map is at least $\sqrt{b}$, the total mass moving by such bad maps is at most $O(\frac{EMD}{\sqrt{b}})$. Therefore,
    $$\left. \text{diff}_p - \text{diff}_{\hat{p}}\right|_{\mathcal{B}^1_1} = \sum_{\hat{K}+\sqrt{b}< t \leq K}(\hat{K}+\sqrt{b}+b-t)\left.(p_t-\hat{p}_t) \right|_{\mathcal{B}^1_1} \leq b \cdot  \left. \sum_{ \hat{K}+\sqrt{b}<t \leq K} \epsilon_t \right|_{\mathcal{B}^1_1} \leq O(\sqrt{b}\cdot EMD),$$
    where $\left. \epsilon_t \right|_{\mathcal{B}^1_1} = \left. p_t - \hat{p}_t \right|_{\mathcal{B}^1_1}$ is the mass change by all such bad maps in $\mathcal{B}^1_1$ at $t$ for any point $\hat{K}+\sqrt{b}<t \leq K$.

    Then, let us consider the set of bad maps $\mathcal{B}^1_2$ from $(\hat{K},\hat{K}+\sqrt{b}]$ to $(\hat{K}+\sqrt{b},K]$. We similarly define $\left. \epsilon_t \right|_{\mathcal{B}^1_2} = \left. p_t - \hat{p}_t \right|_{\mathcal{B}^1_2}$ as the mass change by all such bad maps in $\mathcal{B}^1_2$ at~$t$ for any point $\hat{K}+\sqrt{b}<t \leq K$.

    Since we are considering bad maps of the form $\mathcal{B}^1_2$, we know that $\left. \sum_{ \hat{K}+\sqrt{b}<t \leq K} \epsilon_t \right|_{\mathcal{B}^1_2}$ is at most $\sum_{\hat{K}<t \leq \hat{K}+\sqrt{b}}\hat{p}_t$. Therefore, 
    $$\left. \text{diff}_p - \text{diff}_{\hat{p}}\right|_{\mathcal{B}^1_2} = \sum_{\hat{K}+\sqrt{b}< t \leq K}(\hat{K}+\sqrt{b}+b-t)\left. \epsilon_t \right|_{\mathcal{B}^1_2} \leq b \cdot  \left. \sum_{ \hat{K}+\sqrt{b}<t \leq K} \epsilon_t \right|_{\mathcal{B}^1_2} \leq b \sum_{\hat{K}<t \leq \hat{K}+\sqrt{b}}\hat{p}_t.$$
    Again, this is bounded by $O(\sqrt{b})$ by taking $T=\sqrt{b}$ in Claim~\ref{claim:if AK is opt policy then K to K+T cannot have too much mass}.



    For the set of bad maps of the second kind which are from $(K+\sqrt{b},K+b]$ to $(\hat{K}+\sqrt{b},K]$, we denote it as $\mathcal{B}^2_1$. We know that maps in $\mathcal{B}^2_1$ has nothing to do with $t<\hat{K}+\sqrt{b}$. Moreover, since the moving distance for such map is at least $\sqrt{b}$, the total mass moving by such bad maps is at most $O(\frac{EMD}{\sqrt{b}})$. Therefore,
    \begin{align*}
    \left. \text{diff}_p - \text{diff}_{\hat{p}}\right|_{\mathcal{B}^2_1} &= \sum_{\hat{K} + \sqrt{b}< t \leq K}(K+b-t)\left.(p_t-\hat{p}_t) \right|_{\mathcal{B}^2_1} \\
    &\leq \sum_{\hat{K} + \sqrt{b}< t \leq K}(K-t)\left.\epsilon_t \right|_{\mathcal{B}^2_1} + \sum_{\hat{K} + \sqrt{b}< t \leq K}b\left.\epsilon_t \right|_{\mathcal{B}^2_1} \leq O(\sqrt{b}\cdot EMD)
    \end{align*}
    where $\left. \epsilon_t \right|_{\mathcal{B}^2_1} = \left. p_t - \hat{p}_t \right|_{\mathcal{B}^2_1}$ is the mass change by all such bad maps in $\mathcal{B}^2_1$ at $t$ for any point $\hat{K} + \sqrt{b}< t \leq K$. The last inequality holds since $\sum_{\hat{K} + \sqrt{b}< t \leq K}(K-t)\left.\epsilon_t \right|_{\mathcal{B}^2_1}$ is bounded by $EMD(\hat{p},p)$ and $\sum_{\hat{K} + \sqrt{b}< t \leq K}b\left.\epsilon_t \right|_{\mathcal{B}^2_1} \leq O(b \cdot \frac{EMD}{\sqrt{b}}) = O(\sqrt{b} \cdot EMD)$.

    Then, let us consider the set of bad maps $\mathcal{B}^2_2$ from $(K,K+\sqrt{b}]$ to $(\hat{K}+\sqrt{b},K]$. This is a bit complicated. First, recall that in $\hat{p}$, $A_{\hat{K}}$ is the optimal policy, so $c((A_{\hat{K}})_{\hat{p}}) \leq c((A_{K+\sqrt{b}})_{\hat{p}})$ (note that in case 1, $\hat{K}<K+\sqrt{b}$). We have
    \begin{align}
    & c((A_{\hat{K}})_{\hat{p}}) \leq c((A_{K+\sqrt{b}})_{\hat{p}}) \notag \\
    \iff & \sum_{t \leq  K+\sqrt{b}} (K+\sqrt{b}+b-t)\hat{p}_t \leq (K+\sqrt{b}-\hat{K}) +  \sum_{t \leq \hat{K}}(\hat{K}+b -t )\hat{p}_t \notag \\
    \iff & \sum_{K<t \leq  K+\sqrt{b}} (K+\sqrt{b}+b-t)\hat{p}_t \leq (K+\sqrt{b}-\hat{K}) +  \sum_{t \leq \hat{K}}(\hat{K}+b -t )\hat{p}_t - \sum_{t \leq  K} (K+\sqrt{b}+b-t)\hat{p}_t. \quad \label{eq:ast} \tag{\(\ast\)}
    \end{align}
    We similarly define $\left. \epsilon_t \right|_{\mathcal{B}^2_2} = \left. p_t - \hat{p}_t \right|_{\mathcal{B}^2_2}$ as the mass change by all such bad maps in $\mathcal{B}^2_2$ at~$t$ for any point $\hat{K}+\sqrt{b}<t \leq K$.

    Since we are considering bad maps of the form $\mathcal{B}^2_2$, we know that $\left. \sum_{ \hat{K}+\sqrt{b}<t \leq K} \epsilon_t \right|_{\mathcal{B}^2_2}$ is at most $\sum_{K<t \leq K+\sqrt{b}}\hat{p}_t$. Therefore, 
    $$\left. \text{diff}_p - \text{diff}_{\hat{p}}\right|_{\mathcal{B}^2_2} = \sum_{\hat{K} + \sqrt{b}< t \leq K}(K+b-t)\left.\epsilon_t \right|_{\mathcal{B}^2_2} \leq \sum_{\hat{K} + \sqrt{b}< t \leq K}(K-t)\left.\epsilon_t \right|_{\mathcal{B}^2_2} + \sum_{\hat{K} + \sqrt{b}< t \leq K}b\left.\epsilon_t \right|_{\mathcal{B}^2_2}.$$
    Note that term $\sum_{\hat{K} + \sqrt{b}< t \leq K}(K-t)\left.\epsilon_t \right|_{\mathcal{B}^2_2}$ is bounded by $EMD(\hat{p},p)$ and $\sum_{\hat{K} + \sqrt{b}< t \leq K}b\left.\epsilon_t \right|_{\mathcal{B}^2_2}$ is bounded by $b\sum_{K<t \leq K+\sqrt{b}}\hat{p}_t$. Thus, to show that $\text{diff}_p - \text{diff}_{\hat{p}} \leq O(\sqrt{b} \cdot \max(EMD,1)) - \text{diff}_{\hat{p}}$, it suffices to show that $b\sum_{K<t \leq K+\sqrt{b}}\hat{p}_t \leq O(\sqrt{b}) - \text{diff}_{\hat{p}}$.

    Note that $b\sum_{K<t \leq K+\sqrt{b}}\hat{p}_t$ is less than the left hand side of~\eqref{eq:ast}, hence it suffices to show that
    \[(K+\sqrt{b}-\hat{K}) +  \sum_{t \leq \hat{K}}(\hat{K}+b -t )\hat{p}_t - \sum_{t \leq  K} (K+\sqrt{b}+b-t)\hat{p}_t \leq O(\sqrt{b}) - \text{diff}_{\hat{p}}.\]
    By writing $\text{diff}_{\hat{p}}=\sum_{t \leq \hat{K}+\sqrt{b}}(K-(\hat{K}+\sqrt{b}))\hat{p}_t + \sum_{\hat{K}+\sqrt{b}< t \leq K}(K+b-t)\hat{p}_t - (K-(\hat{K}+\sqrt{b}))$, it suffices to show that (rearranging to make both sides containing positive terms only)
    \[\sum_{t \leq \hat{K}}(\hat{K}+b -t )\hat{p}_t + \sum_{t \leq \hat{K}+\sqrt{b}}(K-(\hat{K}+\sqrt{b}))\hat{p}_t + \sum_{\hat{K}+\sqrt{b}< t \leq K}(K+b-t)\hat{p}_t \leq O(\sqrt{b}) + \sum_{t \leq  K} (K+\sqrt{b}+b-t)\hat{p}_t. \quad \label{eq:chap} \tag{\(\S\)} \]
    Now, compare both sides of~\eqref{eq:chap} for all $t$.

    For $t \leq \hat{K}$, $LHS$ of~\eqref{eq:chap} is $(\hat{K}+b-t+K-(\hat{K}+\sqrt{b}))\hat{p}_t = (b+K-t-\sqrt{b})\hat{p}_t$. $RHS$ of~\eqref{eq:chap} is $(K+\sqrt{b}+b-t)\hat{p}_t$. So, $LHS \leq RHS$.

    For $\hat{K}<t \leq \hat{K}+\sqrt{b}$, $LHS$ of~\eqref{eq:chap} is $(K-(\hat{K}+\sqrt{b}))\hat{p}_t$. $RHS$ of~\eqref{eq:chap} is $(K+\sqrt{b}+b-t)\hat{p}_t$. So, $LHS \leq RHS$ since $t \leq \hat{K}+\sqrt{b}$.

    For $\hat{K}+\sqrt{b}<t \leq K$, $LHS$ of~\eqref{eq:chap} is $(K+b-t)\hat{p}_t$. $RHS$ of~\eqref{eq:chap} is $(K+\sqrt{b}+b-t)\hat{p}_t$. So, $LHS \leq RHS$.

    Therefore,~\eqref{eq:chap} holds. We proved that $\text{diff}_p \leq O(\sqrt{b} \cdot \max(EMD,1))$ and finish the proof of Case 1.

    \textbf{Case 2.} When $K < \hat{K} + \sqrt{b}$. 

    We will briefly give the short proof here as the full procedure is similar to Case 1.

    In this case, by some simple calculations, we have 
    \begin{align*}
    \text{diff}_{\hat{p}} =& \sum_{t> \hat{K}+\sqrt{b}}(\hat{K}+\sqrt{b}-K)\hat{p}_t - \sum_{K < t \leq \hat{K}+\sqrt{b}}(K+b-t)\hat{p}_t\\
    =& (\hat{K}+\sqrt{b}-K) - \sum_{t \leq K}(\hat{K}+\sqrt{b}-K)\hat{p}_t - \sum_{K<t \leq \hat{K}+\sqrt{b}}(\hat{K}+\sqrt{b}+b-t)\hat{p}_t.
    \end{align*}
    By examining different kind of maps, one can check that there are two kinds of bad maps: either from $(\hat{K},\hat{K}+\sqrt{b}]$ to $[1,K]$ (bad maps of the third kind $\mathcal{B}^3$), or from $(\hat{K},\hat{K}+\sqrt{b}]$ to $(\hat{K}+\sqrt{b},\hat{K}+\sqrt{b}+b]$ (bad maps of the fourth kind $\mathcal{B}^4$). We will analyze those bad maps of the third kind using the first expression of $\text{diff}_{\hat{p}}$ and analyze those bad maps of the fourth kind using the second expression in the above formula.

    As always, we can define $\mathcal{B}^3_1$ and $\mathcal{B}^4_1$ respectively for those cases since the moving distance is at least $\sqrt{b}$ by these bad maps. We know that the total mass moving by such bad maps is at most $O(\frac{EMD}{\sqrt{b}})$, thus we have the $O(\sqrt{b} \cdot EMD)$ bound. 

    For those maps with a moving distance less than $\sqrt{b}$ by bad maps, we usually denote them as $\mathcal{B}^3_2$ and $\mathcal{B}^4_2$ in previous arguments. When writing down the expression $\text{diff}_p - \text{diff}_{\hat{p}}$ caused by $\mathcal{B}^3_2$ and $\mathcal{B}^4_2$, bad maps of the form $\mathcal{B}^4_2$ lie in the easier case. One can again use Claim~\ref{claim:if AK is opt policy then K to K+T cannot have too much mass} to get the required $O(\sqrt{b} \cdot EMD)$ bound, by taking $T=\sqrt{b}$. As for bad maps in $\mathcal{B}^3_2$, we consider $c((A_{\hat{K}})_{\hat{p}}) \leq c((A_{K+\sqrt{b}})_{\hat{p}})$. Then~\eqref{eq:ast} holds in exactly the same way. Then, using~\eqref{eq:ast} and $b\sum_{K<t \leq K+\sqrt{b}} \hat{p}_t$ as a bridge, we can write down the similar form as~\eqref{eq:chap}. Finally, we can conclude that $\text{diff}_p - \text{diff}_{\hat{p}} \leq O(\sqrt{b} \cdot \max(EMD,1)) -\text{diff}_{\hat{p}}$. This finishes the proof of Case 2 (hence the lemma).
\end{proof}


\begin{proof}[Proof of Theorem~\ref{thm:main}]
  Theorem~\ref{thm:main} follows directly from the combination of Lemma~\ref{lm:phat Ainf and p AK with EMD O1} to Lemma~\ref{lm:phat is AKhat and p is AK root b loss}.
\end{proof}
